
\section*{摘\qquad 要}
\addcontentsline{toc}{section}{摘\qquad 要}

随着机场地面运行智能化水平的提升，飞机牵引与对接作业对自动化、精确性和安全性提出了更高要求。针对传统人工牵引方式依赖经验、作业效率受限及安全风险较高等问题，本文以自动牵引车与飞机对接过程为研究对象，基于ROS与Gazebo平台构建联合仿真系统。

本文首先对机场牵引作业场景进行抽象，建立双桥转向牵引车运动学模型，并分析转向角、速度与轨迹曲率之间的约束关系。在此基础上，设计飞机辅助感知与相对位姿误差控制方法，通过目标身份确认、相对距离与姿态误差获取，实现牵引车的自动接近、姿态调整和停车判定。随后，基于Gazebo搭建牵引车、飞机目标和机场作业环境，并通过ROS节点完成车辆控制、传感器数据获取、任务模式切换和系统通信。

围绕模型生成、转向桥回中、轮速驱动、目标接近、对接、推出和撤离等环节，本文进行了分阶段仿真测试。测试结果显示，牵引车模型能够在ROS-Gazebo环境中完成低速行驶、前后桥转向、终端停车及任务状态切换等动作；在目标位置和路径发生变化时，系统也能依据误差反馈进行相应调整。该平台为后续样机结构完善、控制参数整定以及多传感器方案迭代提供了可复现实验基础。

\section*{关键词}
自动牵引车；飞机对接；ROS；Gazebo；路径规划；联合仿真
